%%%% ABSTRACT
%%
%% Versão do resumo para idioma de divulgação internacional.

\begin{abstractutfpr}%% Ambiente abstractutfpr
Modern numerical computing and machine learning applications require both mature libraries and the ability to run concurrently, reliably, and with fault tolerance. The Python ecosystem addresses the first requirement through well-established tools, while the Erlang/Elixir ecosystem stands out for the architectural properties of the BEAM VM, which favor distributed and highly available systems. Despite initiatives such as Numerical Elixir, the maturity gap still limits the adoption of Elixir in scenarios that depend on those tools. In this context, this work investigates and develops interoperability mechanisms between Numerical Elixir and Python, resulting in the nx\_py library, which enables tensor transfer between the two environments with and without memory copies. The approach starts from an analysis of communication mechanisms between Elixir and external languages, with emphasis on Native Implemented Functions and the Pythonx library, which allows embedding a CPython interpreter within the BEAM VM. nx\_py lets an Nx tensor be used directly in Python code and decides on its own how to transfer the data: with a copy, when the tensor is stored on Nx.BinaryBackend, and without one, when it is stored on EXLA, in which case Elixir and Python read and write the same region of memory. Empirical validation, conducted through resident memory usage measurements, confirms the elimination of copies in the zero-copy path. It is concluded that the proposed interoperability expands the capabilities of the Elixir ecosystem in this domain while preserving the concurrency and robustness properties of the BEAM VM.
\end{abstractutfpr}
